\section{IL + RL 混合范式：为什么纯 RL 不够}

讲者在 21 分钟段明确提出：\textit{``I strongly believe the answer is no.''}\footnote{来源：\href{https://www.youtube.com/watch?v=ntjOxjZMaac}{Lecture 03 字幕}，区间 00:21:31--00:21:35。} 这里的问题是``纯 RL 能否走完全程''。其论证不是抽象哲学，而是具体失效案例。

\begin{importantbox}{混合训练的工程逻辑}
\begin{enumerate}
    \item 先用 imitation learning 注入人类先验，获得可用但未最优的初始化策略。
    \item 再用 reinforcement learning 在真实反馈上超越演示分布。
    \item 通过 population 对战和对抗样本，修复策略盲区。
\end{enumerate}
\end{importantbox}

\begin{figure}[H]
\centering
\includegraphics[width=0.95\textwidth]{diagrams/il-rl-loop.png}
\caption{讲义整理图：Imitation Learning 与 Reinforcement Learning 的闭环}
\label{fig:il-rl-loop}
\end{figure}

\begin{knowledgebox}{读图：混合训练不是先 IL、后 RL 的一次性接力}
演示提供稳定起点，rollout 暴露长程行为，结果与约束产生学习信号，RL 更新策略，cross-play 再把失败轨迹回流为新数据。循环中任何一层都可能成为瓶颈：演示覆盖不足会限制起点，奖励不完整会鼓励投机，评测对手单一会让漏洞长期隐藏。
\end{knowledgebox}

字幕在 21:06--22:15 这一段给出一个典型现象：纯 RL 能学出能赢的策略，但可能是``unsatisfactory''的赢法，比如把工人全拉去攻击（worker rush/all-in 变体），以 exploit 方式击败特定对手，却不具备广泛可迁移性。

\begin{table}[H]
\centering
\caption{三类训练路径的对比}
\begin{tabular}{p{2.2cm}p{4.0cm}p{4.0cm}p{3.4cm}}
\toprule
路径 & 优势 & 主要风险 & 适用阶段 \\
\midrule
纯 IL & 稳定、样本效率高、可控性强 & 上限受限于演示质量，难以超越专家 & 冷启动与行为边界塑形 \\
纯 RL & 可发现新策略，长期可超人 & 容易投机取巧、奖励黑客、收敛脆弱 & 闭环仿真充分且奖励清晰场景 \\
IL + RL & 吸收先验后再探索，兼顾性能与可控性 & 系统复杂度高，需精细调参与评测 & 开放任务、需兼顾安全与性能 \\
\bottomrule
\end{tabular}
\end{table}

\begin{warningbox}{不要把``赢一局''当作``会这个任务''}
在多策略环境里，局部 exploit 成功可能只是偶然匹配，不代表策略掌握。若缺少跨对手分布验证，线上部署后会快速暴露脆弱点。
\end{warningbox}

把这个结论迁移到 LLM Agent：仅依赖 preference optimization 或少量 reward model 信号，同样可能出现``会做 benchmark，不会做真实任务''。因此课程强调 post-training 的核心不是把平均分再抬高一点，而是把失败模式系统性地压缩。

\subsection{本章小结}
IL + RL 的价值在于把``人类先验''和``环境反馈''联合起来。纯 RL 或纯 IL 都能在局部任务表现良好，但都难以单独支撑开放世界 Agent 的长期鲁棒性。

